%%%%%%%%%%%%%%%%%%%%%%%%%%%%%%%%%%%%%%%%%%%%%%%%%%%%%%%%%%
\section{Numerical Study}\label{sec:numerical}
%%%%%%%%%%%%%%%%%%%%%%%%%%%%%%%%%%%%%%%%%%%%%%%%%%%%%%%%%%

The previous sections characterize the equilibrium in closed form and
derive thresholds that determine the platform's compensation rate,
algorithmic promotion decision, and the creator's effort. This section
shows numerically how these equilibrium objects move with the main
model parameters. The goal is: (i) to illustrate the
comparative statics implied by the theory and (ii) to map the regions in
which the platform and the creator prefer each monetization base.

Section~\ref{subsec:numerics-design} describes the setup and parameter
choices. Section~\ref{subsec:numerics-Q} varies AI baseline quality.
Section~\ref{subsec:numerics-alpha} varies AI learning efficiency.
Section~\ref{subsec:numerics-delta} varies algorithmic influence.
Section~\ref{subsec:numerics-regions} maps the regions in which each
party prefers view-based or engagement-based monetization.

%%%%%%%%%%%%%%%%%%%%%%%%%%%%%%%%%%%%%%%%%%%%%%%%%%%%%%%%%%
\subsection{Setup}
\label{subsec:numerics-design}
%%%%%%%%%%%%%%%%%%%%%%%%%%%%%%%%%%%%%%%%%%%%%%%%%%%%%%%%%%

For each parameter vector $\theta=(Q,\alpha,\delta,t,k)$, we compute
the equilibrium under both monetization bases. The platform's problem
admits candidate solutions in the uncovered and covered demand
regimes. For each candidate, we check the regime-consistency
conditions and keep only those that satisfy them. Among the surviving
candidates, we select the one that gives the highest platform payoff.
This procedure ensures that the reported equilibrium is consistent
with both the platform's optimization problem and the demand regime
it induces.

We track four equilibrium outcomes:
\begin{itemize}
    \item platform utility $u_P$,
    \item creator utility $u_H$,
    \item total views $\mathrm{TV}=D_A+D_H$, and
    \item total engagement $\mathrm{TE}=q_A D_A+q_H D_H$.
\end{itemize}
We report each outcome under view-based monetization (superscript $V$)
and under engagement-based monetization (superscript $E$).

Table~\ref{tab:baseline-parameters} reports the baseline parameter
values and the ranges used in the experiments. The baseline values
satisfy Assumptions~\ref{ass:nonsat} and~\ref{ass:concavity-E} and lie
in an economically interesting region: AI baseline quality is
moderate, so neither side of the market trivially dominates; the AI's
learning rate is positive but well below one, so human effort matters
without fully determining AI quality; and algorithmic influence is
intermediate, so the platform's promotion lever is meaningful but
bounded. When we vary one parameter, we hold the others at their
baseline values.

\begin{table}[ht]
\centering
\caption{Baseline parameters and ranges used in the numerical study}
\label{tab:baseline-parameters}
\begin{tabular}{lccc}
\toprule
Parameter & Description & Baseline & Range \\
\midrule
$Q$      & AI baseline quality           & $0.30$ & $[0.05,1.20]$ \\
$\alpha$ & AI learning efficiency        & $0.30$ & $[0.10,0.60]$ \\
$\delta$ & Algorithmic influence         & $0.40$ & $[0.10,0.50]$ \\
$t$      & Mismatch cost                 & $2.50$ & - \\
$k$      & Algorithmic-neutrality cost   & $1.00$ & - \\
\bottomrule
\end{tabular}
\end{table}

We run four sets of experiments. The first three vary one parameter
at a time: $Q$, then $\alpha$, then $\delta$. The fourth constructs
two-dimensional maps showing which monetization base each party
prefers across pairs of parameters.

%%%%%%%%%%%%%%%%%%%%%%%%%%%%%%%%%%%%%%%%%%%%%%%%%%%%%%%%%%
\subsection{The Effect of AI Baseline Quality}
\label{subsec:numerics-Q}
%%%%%%%%%%%%%%%%%%%%%%%%%%%%%%%%%%%%%%%%%%%%%%%%%%%%%%%%%%

AI baseline quality $Q$ enters the model in two ways. It shifts the
AI's quality through $q_A=\alpha q_H+Q$, and it shifts the demand for
AI content directly. Raising $Q$ therefore makes AI content more
attractive on its own and weakens the platform's reliance on the human
creator. We expect the platform to lower the creator's compensation,
shift promotion toward AI, and let creator effort fall. The
interesting question is whether these forces play out the same way
under both monetization bases.

Figure~\ref{fig:Q-comparative-statics} plots four equilibrium objects
as $Q$ varies, holding $(\alpha,\delta,t,k)$ at their baseline values: platform
utility $u_P$, creator utility $u_H$, total views $\mathrm{TV}$, and
total engagement $\mathrm{TE}$. Each panel shows two lines, one for
each monetization base.

\begin{figure}[ht]
\centering
\textcolor{red}{[Insert six-panel figure: $\beta_H^*$, $q_H^*$, $u_P$,
$u_H$, $\mathrm{TV}$, $\mathrm{TE}$ as functions of $Q$. One line per
monetization base.]}
\caption{Equilibrium outcomes as AI baseline quality $Q$ varies.}
\label{fig:Q-comparative-statics}
\end{figure}

\textcolor{red}{[To fill in: At what value of $Q$ does $u_P^E$ overtake
$u_P^V$? Does $u_H$ fall in $Q$ under view-based monetization? Does
$\beta_H^*$ move toward AI content faster under view-based
monetization? At what $Q$ does the ranking between the contracts
flip?]}

%%%%%%%%%%%%%%%%%%%%%%%%%%%%%%%%%%%%%%%%%%%%%%%%%%%%%%%%%%
\subsection{The Effect of AI Learning Efficiency}
\label{subsec:numerics-alpha}
%%%%%%%%%%%%%%%%%%%%%%%%%%%%%%%%%%%%%%%%%%%%%%%%%%%%%%%%%%

AI learning efficiency $\alpha$ measures how much the AI's quality
improves with the creator's effort. When $\alpha$ is small, the
creator's effort barely affects AI quality, and the platform's only
reason to compensate or promote her is direct human demand. When
$\alpha$ is large, the creator becomes valuable through a second
channel: every unit of her effort also raises the AI's quality. We
expect the platform to promote the creator more as $\alpha$ increases, and we expect this effect to be stronger under
engagement-based monetization, since engagement responds directly to
quality.

Figure~\ref{fig:alpha-comparative-statics} plots the same four
equilibrium objects as functions of $\alpha$, holding $(Q,\delta,t,k)$
at their baseline values.

\begin{figure}[ht]
\centering
\textcolor{red}{[Insert six-panel figure: $\beta_H^*$, $q_H^*$, $u_P$,
$u_H$, $\mathrm{TV}$, $\mathrm{TE}$ as functions of $\alpha$. One line
per monetization base.]}
\caption{Equilibrium outcomes as AI learning efficiency $\alpha$
varies.}
\label{fig:alpha-comparative-statics}
\end{figure}

\textcolor{red}{[To fill in: Does $\beta_H^*$ fall in $\alpha$ under
view-based monetization and rise under engagement-based monetization?
Does creator utility rise under both contracts, or only under
engagement-based monetization? Does total engagement rise faster under
engagement-based monetization?]}



%%%%%%%%%%%%%%%%%%%%%%%%%%%%%%%%%%%%%%%%%%%%%%%%%%%%%%%%%%
\subsection{When Do the Platform and the Creator Agree?}
\label{subsec:numerics-regions}
%%%%%%%%%%%%%%%%%%%%%%%%%%%%%%%%%%%%%%%%%%%%%%%%%%%%%%%%%%

We ask when the platform and the creator agree on the
monetization base and when they disagree. For each parameter vector,
we compute the equilibrium under both bases and check the signs of
$u_P^E-u_P^V$ and $u_H^E-u_H^V$. This produces four regions:

\begin{itemize}
    \item \textbf{Region I:} both parties prefer view-based monetization.
    \item \textbf{Region II:} both parties prefer engagement-based monetization.
    \item \textbf{Region III:} the platform prefers engagement-based monetization, the creator prefers view-based.
    \item \textbf{Region IV:} the creator prefers engagement-based monetization, the platform prefers view-based.
\end{itemize}

Regions III and IV are the disagreement regions. They identify
parameter combinations under which the choice of monetization base is
a source of conflict between the two parties.

We first build the preference map over $(Q,\alpha)$, which connects
directly to the thresholds in Theorems~\ref{lem:creator-choice}
and~\ref{lem:platform-choice}. The theory predicts that view-based
monetization is more attractive when AI baseline quality is low,
while engagement-based monetization becomes more attractive as $Q$ and
$\alpha$ rise.

\begin{figure}[ht]
\centering
\textcolor{red}{[Insert heatmap over $(Q,\alpha)$ with the four
preference regions in four colors. Overlay the theoretical thresholds
$\bar Q$ and $Q^*$.]}
\caption{Platform and creator monetization preferences over AI
baseline quality and AI learning efficiency.}
\label{fig:preference-map-Q-alpha}
\end{figure}

\textcolor{red}{[To fill in: How large is each region? Is there a
substantial Region III (platform wants engagement, creator wants
views)? How well do the numerical region boundaries match the
theoretical thresholds $\bar Q$ and $Q^*$?]}

We then repeat the exercise over $(Q,\delta)$ and $(\alpha,\delta)$ to
ask how algorithmic influence reshapes the preference regions.

\begin{figure}[ht]
\centering
\textcolor{red}{[Insert two heatmaps over $(Q,\delta)$ and
$(\alpha,\delta)$ showing the four preference regions.]}
\caption{Platform and creator monetization preferences as algorithmic
influence varies.}
\label{fig:preference-map-delta}
\end{figure}

\textcolor{red}{[To fill in: Does the disagreement region grow with
$\delta$? Does higher $\delta$ make view-based monetization more
attractive to the platform when $Q$ is high? Does engagement-based
monetization dominate when both $\alpha$ and $\delta$ are high?]}
